\vpage{1}{Compute the representation first}{THREE VIEWS OF ONE COMPUTATION}{Reading source: main text, Sections 2--3; sixth-block supplement, p. 1.}
A dense array, a factored operator, and a compiled execution plan are different objects. They may denote the same linear map, but they need not store the same data, return the same output, or incur the same costs. This appendix makes those distinctions visible through calculations rather than a pipeline of boxes.
\vfig{three_views}
\head{Three models to carry through the paper}
The first is a six-dimensional operator $A=D+uv^T$. Every entry of $A$ is nonzero, yet three $2\times2$ diagonal blocks and two thin vectors suffice to apply it exactly. Pages \pagerefV{3}, \pagerefV{9}, \pagerefV{10}, and \pagerefV{27} reuse this operator or its factors.

The second starts with redundant factors $L$ and $R$. Actual range coordinates reduce a $4\times3$ product to a $2\times3$ core. Pages \pagerefV{11}--\pagerefV{12} connect rank, singular tails, and the geometric reason a small-core optimum is also an ambient optimum.

The third is a seven-node candidate tree with exactly five prunings. Its budget optimum is not selected by any Lagrangian penalty. Pages \pagerefV{21}--\pagerefV{23} let readers inspect every feasible choice and every changed subproblem.
\head{How to use the appendix}
A source line beneath each title links to the relevant preserved manuscript or supplement page in the combined edition. Exact checks use rational, Gaussian-rational, symbolic, and finite-enumeration arithmetic. Smooth drawing coordinates are illustrations of stated formulas, not additional proof evidence.
\takeaway{An exact representation can remove work before a device is chosen. It does not by itself supply a fast kernel, a good approximation, or a fair timing comparison.}

\vpage{2}{The objects and their contracts}{A READER'S ATLAS}{Main text: Sections 2--3, 8--9, and 11.}
The same letter can refer to a mathematical operator or to a particular stored representation. Keep those levels separate when reading an algorithm or a benchmark row.

\begin{tabularx}{\linewidth}{@{}p{.18\linewidth}p{.27\linewidth}Y@{}}\toprule
Symbol & Object & Question to ask \\\midrule
$P_\alpha$ & Coordinate projector & Which input or output coordinates does it select?\\
$A_{\beta\alpha}$ & A block map & What are its source and target dimensions?\\
$E$ & Active block addresses & Is this exact support, or only a structural upper bound?\\
$D+UV^*$ & Factored operator & Is the desired output a vector, factors, or a dense array?\\
$Q_L,G,Q_R$ & Range factors and core & Are the outer columns orthonormal, with exact reconstruction?\\
$E_j$ & An error operator & Is it a factor perturbation or an additive error after a product?\\
$s_e,\epsilon_e$ & Sensitivity and tolerance & Is the budget input-dependent or uniform over unit inputs?\\
$(\kappa_v,\epsilon_v,\sigma_v)$ & Candidate-tree data & Do costs, errors, and split overheads really add?\\
$F_v(b)$ & A budget table & Does $b$ mean exact weight, or an upper bound?\\
$R$ & Reuse count & Which one-time costs have been included?\\\bottomrule
\end{tabularx}
\head{Conventions used in the examples}
Matrices act on column vectors. A block $A_{\beta\alpha}$ goes from source $\alpha$ to target $\beta$. Products act from right to left. Real examples use transpose; complex examples use adjoint unless a vectorization identity explicitly requires transpose. Column-major vectorization stacks the columns of a matrix.

An arithmetic count follows the paper's modeled multiply/add convention. It is not a count of executed machine instructions. The tiny scheduling and amortization examples use declared synthetic units, not measured GPU times. The original benchmark corpus remains historical data in the preserved source tree.
\head{Evidence levels}
The original theorem ledger records its inherited finite formal scope. The new exact checks validate these particular examples. A plot samples a formula. A historical timing describes its recorded case and machine. None of these evidence types can silently substitute for another.
\takeaway{Before interpreting a number, identify its object, norm, output contract, and timing scope. Most apparent contradictions disappear once those four fields are explicit.}

\vpage{3}{A dense-looking matrix, one scalar coupling}{THE RUNNING SIX-DIMENSIONAL OPERATOR}{Main text: Theorems 2.3 and 3.4, pp. 4 and 7.}
Let $D$ have three identical blocks $\mat{3&1\\1&3}$ and put $u=v=\ones\in\QQ^6$. Then $A=D+uv^T$ has diagonal entries four, within-block off-diagonal entries two, and every cross-block entry one. All 36 entries are nonzero.
\[
x=(1,-1,2,0,-2,3)^T,\qquad v^Tx=3,\qquad Ax=Dx+3u.
\]
\vfig{factored_application}
\head{What was avoided?}
The global part does not require 36 independent coefficients. It depends on the input through one scalar, $v^Tx$, which is broadcast through $u$. The diagonal part is computed by three separate local matrix products. Nothing is approximated and no channel is discarded.

The sixth output is $7+3=10$, while the fifth is $-3+3=0$. A zero result can arise through cancellation even when the corresponding row contains no zero coefficient. The output vector is therefore not a support map of the input matrix.
\head{What has not been established?}
This computation proves neither that the supplied factors are the cheapest representation nor that six small tasks outperform one dense kernel. For example, the repeated diagonal blocks and all-one vectors admit further special handling. The stated BDLR model deliberately treats the blocks as dense and the factors as stored vectors so that the paper's general formulas apply literally.

The checker constructs the complete matrix independently, multiplies it by $x$, and compares all six entries with the factored route. The same input returns in the product calculation on \pagerefV{27}.
\takeaway{A matrix can have full entrywise support and still admit a small exact representation. Entry sparsity and algebraic compressibility are different properties.}

\vpage{4}{Work removed is not time measured}{ARITHMETIC, STORAGE, AND CAPACITY}{Main text: Theorems 3.1 and 3.4; Corollaries 3.2 and 3.9, pp. 7--9.}
For $q$ right-hand sides, the paper models dense application by $2n^2q$ operations. A block-diagonal-plus-rank-$r$ representation with equal block width $b$ uses
\[
F=2qn(b+2r),\qquad M=n(b+2r)\text{ stored scalars}.
\]
These formulas count the dense block products and two thin products. They follow the manuscript convention rather than adding every scalar accumulation or index operation.
\vfig{work_storage}
For the running example, $n=6,b=2,r=1$. The model reduces storage from 36 to 24 scalars and application work from $72q$ to $48q$. For fixed $b,r$, the ratio $n/(b+2r)$ grows with $n$. If the rank instead grows proportionally with dimension, that apparent advantage need not grow.
\head{Memory is another constraint}
With scalar size $s$, the representation alone requires $sM$ bytes. Inputs, outputs, workspace, allocator overhead, and stored indices are additional. A matrix-only capacity inequality is not a complete device-memory fit test.

For a block band of $p$ packets and half-bandwidth $w$, the exact active-block count is
\[
p(2w+1)-w(w+1).
\]
The subtraction removes overcounted edge diagonals. The checker enumerates all addresses for $1\le p\le6$ and every admissible $w$, including $w=0$ and $w=p-1$.
\takeaway{An arithmetic or storage ratio explains a representation advantage. Kernel launch cost, traffic, utilization, and conversion determine whether that advantage becomes a time advantage.}

\vpage{5}{Four coefficients, four exact frequencies}{CIRCULANT MULTIPLICATION}{Main text: Theorem 4.3, p. 11.}
Take $c=(2,1,0,-1)^T$ and define $C(c)_{jk}=c_{j-k\bmod4}$. The unitary Fourier matrix is $F_{jk}=\mathrm i^{-jk}/2$. Its normalization is visible in the exact diagonalization:
\[
C(c)=F^*\diag(2Fc)F,\qquad 2Fc=(2,2-2\mathrm i,2,2+2\mathrm i)^T.
\]
\vfig{circulant}
\head{Follow an actual input}
For $x=(1,2,-1,0)^T$, dense multiplication gives $(0,6,0,-2)^T$. Transforming $x$, multiplying by the four displayed eigenvalues, and applying $F^*$ gives exactly the same vector. Every entry belongs to $\QQ(\mathrm i)$, so the check needs no approximate trigonometric values.

There is a common normalization trap. With an unnormalized discrete transform, the convolution multiplier has one convention. With the unitary transform used here, the eigenvalue vector is $\sqrt n\,Fc$, not $Fc$. Dropping the factor two reconstructs $C/2$ rather than $C$ in this example. The invalid control deliberately tries that change and rejects it.
\head{Where the fast algorithm enters}
The identity specifies the represented operator. An FFT is an additional factorization that evaluates the transforms efficiently. The algebraic diagonalization does not certify a particular floating-point FFT, data layout, or launch policy. For a tiny four-point example, direct multiplication may be the simpler execution path even though the frequency representation is exact.
\takeaway{Transform coordinates change both the storage pattern and the normalization conventions. Verify the inverse transform and convolution scale before interpreting a frequency-domain product.}

\vpage{6}{Padding is part of the theorem}{TOEPLITZ EMBEDDING WITHOUT WRAPAROUND}{Main text: Corollary 4.4, p. 12.}
Consider
\[
T=\mat{2&-1&5\\3&2&-1\\4&3&2},\qquad x=(1,2,-1)^T.
\]
The diagonal pattern is Toeplitz, not circulant. To use a circular convolution, embed it into dimension six with first column $g=(2,3,4,0,5,-1)^T$ and pad the input with three zeros.
\vfig{toeplitz}
\head{Read the wrapped indices}
For output index $j<3$, an input index $k\le j$ uses the first-column coefficient $c_{j-k}$. An input index $k>j$ wraps to the stored first-row coefficient $r_{k-j}$. The padded input positions contribute nothing. Keeping the first three outputs therefore gives
\[
Tx=(-5,8,8)^T.
\]
The remaining circular outputs are workspace consequences of the embedding, not additional components of $Tx$.
\head{A concrete failure mode}
Using only the first column $(2,3,4)$ as a three-dimensional circulant changes the upper-triangular coefficients. That is a different matrix. The checker compares its output with $Tx$ and rejects the unpadded shortcut.

The middle zero in $g$ is not a discarded Toeplitz coefficient: it separates the lower and reversed upper diagonals in the chosen $2n$ embedding. Other embeddings require their own index identities. It is not sufficient for an array merely to have the right length or to produce a plausible plot.
\takeaway{Zero padding prevents the wrong circular interactions. It is an algebraic part of the operator representation, not an optional performance tweak.}

\vpage{7}{A transpose, not an adjoint}{KRONECKER ACTION AND ARRAY ORDER}{Main text: Theorem 4.5 and Corollary 4.6, p. 12.}
With column-major vectorization, the tensor identity is
\[
(A\otimes B)\operatorname{vec}(X)=\operatorname{vec}(BXA^T).
\]
Here the transpose records index order, not a Hermitian inner product. Use the explicitly complex matrix $A=\mat{1&\mathrm i\\2&-1}$, together with $B=\mat{2&1\\0&3}$ and $X=\mat{1&2\\3&4}$.
\vfig{kronecker}
\head{Check the coordinates}
The stacked input is $(1,3,2,4)^T$. Direct tensor multiplication agrees with stacking the columns of
\[
BXA^T=\mat{5+8\mathrm i&2\\9+12\mathrm i&6}.
\]
Replacing $A^T$ with $A^*$ changes the signs of the imaginary terms. A real-only test would conceal that error, which is why the worked model uses an imaginary coefficient.
\head{Why modewise contraction matters}
For $A$ of order $m$ and $B$ of order $n$, the explicit tensor matrix stores $m^2n^2$ entries. Its factors store $m^2+n^2$. Applying through $BXA^T$ costs $2n^2m+2nm^2$ modeled operations, rather than forming the full tensor matrix first.

A different flattening convention has a different coordinate formula. The statement is not independent of how an array is reshaped. In a software implementation, shape, stride, and flattening order belong to the correctness contract as well as the cost model.
\takeaway{Tensor structure gives exact savings only when the address convention is preserved. A silent transpose, conjugation, or reshape change implements a different map.}

\vpage{8}{See every butterfly value}{WALSH RECURSION IN EIGHT COORDINATES}{Main text: Theorem 4.7, p. 13.}
The unnormalized Walsh matrix satisfies $H_1=[1]$ and
\[
H_{2n}=\mat{H_n&H_n\\H_n&-H_n},\qquad H_nH_n^T=nI.
\]
Apply the three add/subtract stages to $(1,2,3,4,5,6,7,8)^T$. The full intermediate ledger is displayed below.
\vfig{walsh}
\head{Arithmetic can be counted exactly}
Each stage visits eight coordinates, producing one sum or difference per coordinate. Three stages therefore use exactly 24 additions or subtractions. No multiplication by a dense $8\times8$ matrix is necessary.

The result $(36,-4,-8,0,-16,0,0,0)^T$ still reconstructs the input: the inverse is $H_8^T/8$. Because the displayed transform is unnormalized, its squared output norm is eight times the squared input norm. The normalized transform $H_8/\sqrt8$ would instead be orthogonal.
\head{What the zeros mean}
The zero coefficients arise from exact cancellation among the input values. They do not imply that the transform's defining matrix is sparse. The butterfly structure, not the final coefficient support, explains the arithmetic reduction.

The checker verifies the recursion, orthogonality, dense-versus-butterfly output, and operation count for sizes one, two, four, eight, and sixteen. These bounded tests exercise the implementation and boundary cases. The general recursion theorem remains the manuscript's mathematical statement.
\takeaway{Factorization can turn a dense matrix into a short exact arithmetic circuit. Keep its normalization separate from its recursive evaluation rule.}

\vpage{9}{The cross term has somewhere to go}{CLOSED LOW-RANK PRODUCT ALGEBRA}{Main text: Theorem 5.1 and Corollary 5.2, pp. 13--14.}
Let $A=D_A+u_av_a^T$ be the running operator. Take $D_B=\diag(1,2,3,4,5,6)$, $u_b=(1,0,1,0,1,0)^T$, and $v_b=(0,1,0,1,0,1)^T$. Define $B=D_B+u_bv_b^T$.
\[
U_C=[D_Au_b+u_a(v_a^Tu_b)\ \ u_a],\qquad
V_C=[v_b\ \ D_B^Tv_a].
\]
\vfig{bdlr_product}
\head{Why there are only two columns}
Expanding $AB$ gives three correction terms, but two share the right factor $v_b^T$:
\[
D_Au_bv_b^T+u_a(v_a^Tu_b)v_b^T
=\bigl(D_Au_b+u_a(v_a^Tu_b)\bigr)v_b^T.
\]
The remaining term is $u_av_a^TD_B$. Thus $AB=D_AD_B+U_CV_C^T$ with at most two correction directions. Counting the three expanded summands as independent ranks would overestimate this representation.

The scalar $v_a^Tu_b$ is three. Removing the quadratic term would therefore remove $3u_av_b^T$, which is nonzero. The checker compares the complete product to the factor construction and rejects a calculation that omits this contribution.
\head{Before and after recompression}
A rank bound on stored factors is not always the actual matrix rank. Repeated products may introduce dependent columns. Recompression can remove those dependencies exactly, or remove small singular directions approximately under an explicit error budget.
\takeaway{Closure of a representation class is an algebraic statement about all product terms. It is not permission to drop an interaction merely because it is second order in the factors.}

\vpage{10}{A large inverse, a small denominator}{WOODBURY WITH BOTH INVERTIBILITY CHECKS}{Main text: Theorem 5.3 and Corollary 5.4, pp. 14--15.}
For the running operator, each diagonal block has inverse $\frac18\mat{3&-1\\-1&3}$. Hence $D^{-1}u=u/4$ and
\[
S=1+u^TD^{-1}u=\frac52.
\]
The Woodbury identity reduces the entire inverse to
\[
A^{-1}=D^{-1}-D^{-1}uS^{-1}u^TD^{-1}
=D^{-1}-\frac1{40}uu^T.
\]
\vfig{woodbury}
\head{Two hypotheses, not one}
Invertibility of $D$ alone does not suffice. For $D=I_2$, $u=e_1$, and $v=-e_1$, the update is $\diag(0,1)$. The small matrix $S$ is zero, exactly where the correction formula would require an invalid division.

For the running model, the checker multiplies the proposed inverse by $A$ in both orders and obtains the identity. It also checks the determinant relation $\det A=\det D\det S$. These tests verify the actual factors and sign of the correction, not just a numerical solve residual.
\head{Conditioning survives the reduction}
A small $S$ can still be ill-conditioned. Reducing its dimension does not bound $\|S^{-1}\|$. The manuscript's inverse estimate keeps that term explicit, along with the norms of $D^{-1}u$ and $v^TD^{-1}$.
\takeaway{The hidden coupling has not vanished. Its solvability and sensitivity are concentrated in a smaller operator that still has to be invertible and numerically controlled.}

\vpage{11}{Stored width is not intrinsic rank}{ACTUAL RANGE COORDINATES}{Main text: Theorem 5.6, p. 15; sixth-block supplement, p. 2.}
Choose the isometry
\[
Q_L=\mat{3/5&0\\4/5&0\\0&1\\0&0},\quad
R_L=\mat{1&0&1\\0&1&0},\quad
R_R=\mat{3&0&2\\0&2&0\\1&0&-1}.
\]
Set $L=Q_LR_L$ and $R=R_R$. The two equal columns of $L$ do not create a third range direction. In contrast, $R$ has rank three.
\vfig{compact_core}
The actual core is
\[
G=R_LR_R^T=\mat{5&0&0\\0&2&0},\qquad LR^T=Q_LG.
\]
Its dimensions are the two actual input ranks, two and three. This is not inferred by merely discarding a column with a small norm. The range factor $Q_L$ has orthonormal columns and reconstructs $L$ exactly.
\head{The preserved formal boundary}
The sixth release constructs orthonormal bases of the actual column spaces. It then assumes exact ordered thin-SVD data for the core, as the manuscript instructs the reader to take an SVD. It does not assume an optimality or tail-error conclusion. Nor does it certify a floating-point QR or SVD implementation simply because those routines use similar notation.

This example has a diagonal rectangular core, so its singular directions are explicit. No iterative numerical factorization is needed to identify the two nonzero singular values, five and two.
\takeaway{Compactness is a statement about actual range spaces. The number of stored factor columns can exceed the dimensions that matter for recompression.}

\vpage{12}{Why a core optimum is global}{COMPRESSION, KERNEL WITNESSES, AND TAILS}{Reading source: sixth-block supplement, pp. 3--6.}
The product from the previous page is
\[
A_0=LR^T=\mat{3&0&0\\4&0&0\\0&2&0\\0&0&0}.
\]
Keeping only its first singular direction gives a rank-one approximation with spectral and Frobenius error exactly two. Keeping none gives errors five and $\sqrt{29}$; keeping both gives zero.
\vfig{optimality}
\head{Why an outside competitor cannot do better}
For isometric outer factors, compression $C(B)=Q_L^*BQ_R$ cannot increase rank or either matrix norm. Every ambient rank-$k$ competitor therefore supplies an admissible core competitor. The lifted core residual has exactly the core residual norm. These two facts transfer a core optimum to the entire ambient matrix space, not just to competitors already lying in the supplied ranges.

For a $d$-dimensional diagonal core and $0\le k<d$, the spectral lower bound has a concrete witness. A rank-at-most-$k$ competitor has a nonzero kernel vector in the first $k+1$ coordinates. On that vector the diagonal matrix acts with norm at least $\sigma_{k+1}$ times the vector norm. The competitor cannot cancel that direction.

The Frobenius proof uses an orthonormal basis of the competitor's kernel. Its row weights lie in $[0,1]$ and sum to the kernel dimension. Weighting the decreasing squared singular values forces at least the discarded tail. The two norm arguments are distinct; one is not a graph of the other.
\takeaway{Optimality is not inferred from the appearance of the singular-value plot. It follows from compression and a lower bound valid against every rank-bounded competitor.}

\vpage{13}{Errors enter at different places}{FACTOR PERTURBATION VERSUS RECOMPRESSION}{Main text: Theorems 5.8--5.9, p. 17.}
Let the exact scalar factors be $(2,3,4)$ and start from $p_0=1$. If additive errors $(1/10,1/5,1/20)$ enter \emph{after} each partial multiplication, then
\[
\widehat p_j=a_j\widehat p_{j-1}+e_j,\qquad
\widehat p_3-p_3=12e_1+4e_2+e_3=\frac{41}{20}.
\]
\vfig{error_chain}
\head{The same local numbers, a different model}
If instead the factors themselves are perturbed, the computed product is $(4+e_3)(3+e_2)(2+e_1)$. Its error is $402/125$, not $41/20$. The factor-perturbation telescoping formula contains already-perturbed later factors; the additive recompression formula contains the original later factors.

The visual comparison uses the additive model. The first error is multiplied by both remaining factors. The last error is not amplified at all. Summing only the local error magnitudes gives $7/20$, which is smaller than the actual final error and therefore not a valid bound here.
\head{In a matrix chain}
The corresponding formulas replace scalar magnitudes with operator norms and respect multiplication order. A bound can be conservative when contributions cancel, but it must not omit possible amplification. Singular-tail errors from a recompression step become useful only after they are inserted into the correct global error model.
\takeaway{An error tolerance needs an entry point. A bound for perturbing factors and a bound for compressing partial products are not interchangeable.}

\vpage{14}{Two routes can cancel exactly}{SUPPORT IS AN INCLUSION}{Main text: Theorems 2.5 and 6.1; Corollary 6.3, pp. 5 and 18.}
Let $K$ have transfers $0\to1$, $0\to2$, $1\to3$, and $2\to3$, with coefficients $1,1,1,-1$. There are two length-two routes from packet zero to packet three, but their contributions sum to zero.
\vfig{cancellation}
\[
(K^2)_{30}=K_{31}K_{10}+K_{32}K_{20}=1-1=0,
\qquad K^2=0.
\]
The matrix of path counts sees the two possibilities and forgets the signs. It therefore predicts possible support but not actual support. Replacing the inclusion $\supp(AB)\subseteq\supp(A)\circ\supp(B)$ by equality would be wrong even in this four-dimensional model.
\head{What an acyclic graph does guarantee}
If the longest strict directed path has $h$ edges, then no term of $K^{h+1}$ can survive. Thus $K^{h+1}=0$. This is an upper bound on the nilpotence index, not an assertion that $K^h\ne0$. Here the graph has paths of length two, while coefficient cancellation already gives $K^2=0$.

A structurally permitted block may also vanish because a product of rectangular maps is zero before any sum is taken. A support graph cannot distinguish that mechanism from cancellation among separate routes. The actual block coefficients are therefore part of a proof-bearing operator representation.
\takeaway{Use graph reachability to exclude impossible transfers. Use the algebra to determine which permitted transfers actually survive.}

\vpage{15}{A solve has a dependency clock}{TOPOLOGICAL LEVELS WITH EXACT OUTPUTS}{Main text: Theorem 6.5, Corollary 6.6, and Theorem 10.3, pp. 18--19 and 32.}
For seven scalar packets, put two on the diagonal and $-1$ on each declared incoming edge. The edges are $0\to2$, $1\to2$, $1\to3$, $2\to4$, $3\to4$, $3\to5$, $4\to6$, and $5\to6$. Solve $(2I+K)x=\ones$.
\vfig{triangular}
\head{Read the stages as dependencies}
Packets zero and one are ready first. Packets two and three use those values. Packets four and five then become available, followed by packet six. For instance,
\[
x_4=\frac{1+x_2+x_3}{2}=\frac{11}{8},\qquad
x_6=\frac{1+x_4+x_5}{2}=\frac{13}{8}.
\]
All entries in the displayed solution are checked by multiplying the original triangular matrix by the resulting vector.
\head{When is four a minimum?}
There are four vertices on a longest dependency path. Under the paper's unit-duration, unlimited-concurrency model, four stages are both necessary and attainable. If only one execution slot is available, seven tasks require seven slots of work. If task durations differ, the relevant critical path is weighted.

The theorem is not a universal GPU scheduling result. Transfers, synchronization, contention, and nonunit task costs change the model. The exact benefit here is that independent rows in one level need not wait for one another's mathematical results.
\takeaway{A topological schedule separates dependencies from resources. Its optimal depth is conditional on the task-duration and concurrency model that defines a stage.}

\vpage{16}{Elimination creates new interactions}{SCHUR FILL AND ITS CANCELLATION}{Main text: Theorem 6.7 and Proposition 6.8, pp. 19--20.}
Retain the first two coordinates of
\[
A=\mat{2&0&1\\0&2&1\\1&1&2}.
\]
The third coordinate is hidden. Its invertible block is $[2]$, so eliminating it produces
\[
S=A_{RR}-A_{RH}A_{HH}^{-1}A_{HR}
=\mat{3/2&-1/2\\-1/2&3/2}.
\]
\vfig{schur_fill}
\head{The hidden route has a sign}
The retained coordinates had no direct coupling, yet both touch the hidden coordinate. The resulting correction couples them by $-1/2$. This is fill: a new active retained block produced by an eliminated route.

Now change both original retained off-diagonal entries to $+1/2$. The same hidden correction cancels them, leaving the diagonal matrix $\frac32 I$. Thus a structural route can create or remove visible coupling depending on the direct coefficient already present.
\head{A stronger inverse statement}
For an invertible finite packet matrix, inverse support lies in the directed transitive closure. A reachable-coordinate subspace is invariant under the matrix. Injectivity of its restriction implies surjectivity in finite dimension, so the inverse preserves it as well. This proof uses finite-dimensional invertibility, not a convergent geometric series.

The checker verifies both Schur complements and, in each invertible example, equality between $S^{-1}$ and the retained block of $A^{-1}$.
\takeaway{Removing coordinates does not remove their influence. It changes the retained operator, with possible fill, cancellation, and new conditioning constraints.}

\vpage{17}{Invariant does not mean independently addressable}{A STRICT KRYLOV INCLUSION}{Main text: Proposition 7.2 and Theorem 7.3, p. 21.}
Take $A=I_2$, split into the two coordinate packets, and start with $v=(1,1)^T$. Every packet is invariant. Nevertheless every polynomial applies the \emph{same} scalar to both coordinates:
\[
p(A)v=p(1)(1,1)^T,\qquad \dim\mathcal K_m(A,v)=1\quad(m\ge1).
\]
\vfig{krylov}
The direct sum of the two local Krylov spans is all of $\QQ^2$. It permits independent coefficients in the two packets. The global Krylov space does not. This is why the manuscript states an inclusion, not an unconditional equality.
\head{What breaks the linkage?}
With $A=\diag(1,2)$, the vectors $v$ and $Av=(1,2)^T$ are independent. Two polynomial samples now distinguish the packet labels. Alternatively, allowing the packet projectors as extra address operators separates the coordinates even when $A=I_2$. Both changes enlarge the addressing power, but in different ways.
\head{When invariance is only approximate}
If $A=D+K$ with $D$ packet diagonal, then
\[
\|A^j-D^j\|\le jM^{j-1}\|K\|,
\qquad M=\max\{\|A\|,\|D\|\}.
\]
This controls departure from the decoupled iterate. It is not a claim that the dimension or numerical conditioning of a Krylov basis varies continuously. Near dependence can remain difficult even when the operator error is small.
\takeaway{A single global polynomial can lock otherwise invariant packets together. Independence of subspaces is not the same as independence of the allowed actions.}

\vpage{18}{Measure coupling after normalization}{BLOCK JACOBI AS A GEOMETRIC CHANGE}{Main text: Theorem 7.4 and Corollary 7.6, p. 22.}
Let
\[
A=\mat{4&1\\1&1},\quad D=\diag(4,1),\quad
B=D^{-1/2}=\diag(1/2,1).
\]
Then $BAB=I+E=\mat{1&1/2\\1/2&1}$. Its exact eigenvalues are $1/2$ and $3/2$, giving condition number three.
\vfig{preconditioning}
\head{The unit circles of the energy norms}
The left ellipse is $x^TAx=1$. Writing $x=By$ gives the right ellipse $y^T(BAB)y=1$. The drawing samples these explicitly defined curves; the norm and eigenvalue equalities are checked symbolically.

Here $\|E\|_2=1/2$. The manuscript's general inequality
\[
\kappa_2(D^{-1/2}AD^{-1/2})
\le\frac{1+\eta}{1-\eta},\qquad \|E\|_2\le\eta<1,
\]
is attained. The raw off-diagonal entry is one; its meaningful size is one-half only after normalization by the diagonal energy scales.
\head{Keep the convergence input separate}
The condition bound can be inserted into a conjugate-gradient convergence estimate under that method's positive-definite assumptions. The CG estimate is one of the separately tracked classical inputs. The algebraic normalization theorem does not certify every stopping rule, floating-point iteration, or arbitrary preconditioner.
\takeaway{A small raw coupling is not scale-invariant. Normalize against the diagonal blocks before deciding whether the coupling is weak enough for a condition bound.}

\vpage{19}{Twelve small blocks, one coherent error}{DISCOVERY AND PERTURBATION CERTIFICATES}{Main text: Theorem 8.1 and Proposition 9.2, pp. 23 and 26.}
Split $\QQ^4$ into scalar packets and let
\[
A=I+E,\qquad E=\frac1{100}(\ones\ones^T-I).
\]
Dropping every off-diagonal block leaves $I$. Each discarded entry is only $1/100$, but the total error is not bounded by that single-entry size.
\vfig{block_error}
\head{Two norms answer different questions}
There are twelve omitted entries, so $\|E\|_F^2=12/10^4$. The all-one vector is an eigenvector with eigenvalue $3/100$, while the three orthogonal directions have eigenvalue $-1/100$. Hence $\|E\|_2=3/100$.

For orthogonal coordinate blocks, the squared Frobenius norms add exactly. Operator norms instead require a suitable bound or further spectral information. The maximum omitted block norm is generally insufficient.
\head{Input sensitivity is additional information}
The packet perturbation estimate is
\[
\|(\widehat A-A)x\|\le\sum_e\|E_e\|\,\|P_{\alpha(e)}x\|.
\]
It can be much tighter for an input that scarcely occupies an affected source packet. To obtain a uniform operator bound, the input dependence must be bounded uniformly over unit vectors. A certificate for one input is not automatically a certificate for every right-hand side.

These are a posteriori checks for supplied projectors. They do not discover the best packet system or prove that a thresholding heuristic is globally optimal.
\takeaway{A proposed compression needs a residual certificate in the requested norm and input scope. Small individual blocks can align into a larger collective error.}

\vpage{20}{Spend error where it is cheapest}{CONTINUOUS AND CLIPPED PRECISION}{Main text: Theorems 8.4--8.5, p. 24.}
Take sensitivities $s=(1,2,4)$, cost coefficients $c=(1,1,2)$, and budget $\tau=1$. Minimizing $\sum_ec_e\log(1/\epsilon_e)$ subject to $\sum_es_e\epsilon_e=1$ gives
\[
\epsilon_e=\frac{c_e}{4s_e},\qquad
\epsilon=(1/4,1/8,1/8).
\]
\vfig{precision}
\head{The budget is weighted}
The three contributions to the error allowance are $(1/4,1/4,1/2)$. Equal cost coefficients give equal weighted contributions, not equal tolerances. The third edge has twice the cost coefficient of the second, which offsets its doubled sensitivity in this example.

Now impose upper tolerances $m=(1/10,1/2,1/2)$. The first edge must stop at its tighter cap. With multiplier $\lambda=10/3$, the clipped solution is
\[
\epsilon_e=\min\{m_e,c_e/(\lambda s_e)\},\qquad
\epsilon=(1/10,3/20,3/20).
\]
Its weighted contributions are $(1/10,3/10,3/5)$ and still sum to one. The checker verifies the active cap, the multiplier equations, and the budget exactly.
\head{From tolerance to a real data type}
The continuous objective is a model. Actual kernels offer discrete precisions and may use different accumulation types. A tolerance below one does not identify a particular hardware format. The finite option program on the next page is the correct next step once admissible costs and error bounds have been supplied.
\takeaway{Precision allocation is an optimization over certified sensitivities and costs. It does not infer error bounds from a data-type name alone.}

\vpage{21}{An exact budget is not a greedy choice}{THREE EDGES, TWENTY-SEVEN ASSIGNMENTS}{Main text: Theorem 8.7 and Remark 8.8, p. 26.}
Each edge has three allowed $(\text{integer weight},\text{cost})$ options:
\[
\begin{aligned}
1 &: (0,7),(2,3),(4,1),\\
2 &: (0,6),(1,4),(3,1),\\
3 &: (0,9),(2,5),(5,1).
\end{aligned}
\]
There are exactly $3^3=27$ full assignments. The table stores the minimum cost at each \emph{exact} weight; columns zero through eight are displayed.
\vfig{discrete}
\head{Read the final row correctly}
At budget four, minimize across columns zero through four. The optimum costs 14: choose weights $(2,0,2)$ with costs $(3,6,5)$. The zero-weight choice on edge two is its first, more expensive option. A locally cheap choice can consume too much of the shared budget.

The checker compares the recurrence with complete enumeration for every budget from zero through twelve. Infeasible states remain $+\infty$; they are not zero-cost states.
\head{Quantization must be conservative}
Suppose real errors are $3/50$ and $2/25$, and one integer unit represents $1/10$. Rounding down gives two zero weights, even though the true sum is $7/50>1/10$. Rounding up gives a certified upper envelope instead. The budget method is exact for its integers; it cannot repair unsafe integers supplied by a preprocessing step.
\takeaway{The table optimizes the declared option set exactly. Accuracy depends separately on the validity and conservativeness of the option weights.}

\vpage{22}{A penalty can miss the budget optimum}{FIVE PRUNINGS, ONE NONCONVEX FRONTIER}{Main text: Theorem 9.4, Corollary 9.7, and Theorem 9.8, pp. 27--30.}
A root $r$ has children $a,b$, each with two terminal children. Keeping $r$ costs one with error six; splitting it adds overhead one. Keeping either $a$ or $b$ costs three with error two. Splitting $a$ costs four with zero error, and splitting $b$ costs five with zero error.
\vfig{tree_frontier}
The five complete prunings have $(\text{cost},\text{error})$ pairs
\[
(1,6),\ (7,4),\ (9,2),\ (8,2),\ (10,0).
\]
At budget two the unique cheapest pruning is $\{a_0,a_1,b\}$, costing eight. It is Pareto optimal, but never minimizes $\cost+\lambda\err$ for any $\lambda\ge0$.
\head{See the incompatible inequalities}
To beat the root, that pruning would require $8+2\lambda\le1+6\lambda$, hence $\lambda\ge7/4$. To beat the full split, it would require $8+2\lambda\le10$, hence $\lambda\le1$. No multiplier meets both conditions.

The exact budget recurrence retains the point because it enumerates admissible weight allocations. A sweep of Lagrangian parameters sees only supported frontier points. The two compiler formulations are therefore genuinely different, even on this tiny tree.
\takeaway{A penalty parameter is not a universal substitute for an explicit budget. Discrete, nonconvex frontiers can contain useful optima that no penalty selects.}

\vpage{23}{Recompile the changed part, not everything}{LOCAL DEPENDENCIES AND STRICT MARGINS}{Main text: Theorems 9.11--9.12, pp. 30--31.}
Use the same tree with $\lambda=2$. Increase only the leaf cost at $a_0$ from two to three. The exact subtree values change at $a_0$, its parent $a$, and the root $r$. Every value in the $b$ subtree remains untouched.
\vfig{local_update}
\head{Why the update is local}
A dynamic-programming value depends only on the local keep/split data and its children's values. A changed leaf can therefore affect its ancestors but not a disjoint sibling subtree. Recomputing the dependency closure yields the same global optimum as a fresh traversal. In this example, the root value increases from ten to eleven.
\head{Strict gaps protect decisions}
Before the update, the root keep score is 13 and its split score is 10, giving gap three. If each score changes by less than $3/2$, their strict ordering is retained. Allowing exactly $3/2$ in both directions can produce a tie, so the strict inequality in the stability statement matters.

The bound applies to the completed scores after child subproblems have been evaluated. Bounds on raw leaf perturbations must first be propagated to those scores. Shared caches, overlapping subproblems, or global rank constraints can introduce dependencies outside a tree and invalidate this simple locality argument.
\takeaway{Local recompilation is exact because the model has local dependencies. The locality is a mathematical hypothesis, not an automatic property of every compiler optimization.}

\vpage{24}{Batch shapes, account for traffic}{A LAUNCH MODEL IS NOT A BANDWIDTH MODEL}{Main text: Proposition 3.6, Theorem 3.8, and Proposition 10.2, pp. 9 and 31.}
Six independent blocks have the displayed dimensions. When right-hand-side count, scalar type, accumulation type, and kernel contract also agree within each class, they can be grouped into three batches. That reduces the modeled launch count from six to three without changing their arithmetic.
\vfig{batch_shapes}
\head{Independence must be real}
The batching statement assumes compatible backend support and independent block tasks. If several blocks accumulate into one target, the reduction or accumulation policy is additional work. Equal matrix dimensions alone do not remove that dependency. Different accumulation types also define different shape classes under the paper's convention.
\head{Count the data separately}
For three diagonal $2\times2$ blocks, eight-byte scalars, and $q$ right-hand sides, the stated one-pass model moves
\[
B_{\rm stream}=8\cdot3(4+2q+4q)=96+144q\text{ bytes}.
\]
It gives 240 bytes at $q=1$ and 672 bytes at $q=4$. It deliberately assumes no reuse across edges. A cache-aware implementation may have another traffic model, but it must state one.

Given operation capacity $P$ and bandwidth $W$, $T\ge\max(F/P,B/W)$ is a lower bound, not an achieved runtime. Launch overhead, imperfect overlap, and utilization can make the actual time much larger.
\takeaway{Grouping can remove launches while leaving the same products and data movement. Arithmetic, traffic, and scheduling are related constraints, not one interchangeable cost.}

\vpage{25}{The first win comes after the tie}{SETUP AND REUSE IN ONE LINEAR MODEL}{Main text: Theorem 10.4 and Corollaries 10.5--10.6, pp. 32--33.}
Suppose the baseline costs seven units per use. A candidate costs 30 units to set up and two per use. These are synthetic constants chosen to expose the exact threshold:
\[
T_b(R)=7R,\qquad T_c(R)=30+2R.
\]
\vfig{amortization}
At $R=6$, both cost 42. The candidate becomes strictly faster at $R=7$, where it costs 44 instead of 49. The minimum profitable integer reuse count is
\[
\left\lfloor\frac{30}{7-2}\right\rfloor+1=7.
\]
Replacing the floor-plus-one rule with a ceiling would be wrong when the ratio is already an integer and strict improvement is required.
\head{Which costs belong to setup?}
Discovery, compilation, host/device transfer, graph capture, and factor preparation can have different reuse lifetimes. If a matrix changes but its shape does not, a captured schedule might remain reusable while its numerical factors do not. A single setup number is meaningful only after those lifetimes are fixed.

Materialization has the opposite tradeoff: pay once for an explicit representation, then use a potentially faster dense downstream operation. It pays off only when the explicit per-use cost is lower and the reuse count exceeds the corresponding conversion threshold.
\takeaway{A faster steady-state kernel can lose end-to-end. State the reuse regime and count the tie separately from the first strict win.}

\vpage{26}{Same output before comparing times}{A BENCHMARK ROW IS A CONTRACT}{Main text: Proposition 2.9 and Definitions 11.1--11.7, pp. 6 and 33--34.}
Consider four synthetic timing rows for a product. Row B returns factors, while row A returns an explicit matrix. A ratio $12/2=6$ does not compare two implementations of the same returned object. Adding the decoding cost changes the candidate to row C.
\vfig{contract_match}
\head{What a fair match retains}
The input values, dimensions, layout, source precision, correctness reference, returned output, and timing scope must agree. Row D changes the input precision; that can be a useful accuracy tradeoff, but it is not the same case merely because the output has the same shape.

If a downstream task consumes factors directly, row B has a valid compressed-output contract. It should be compared with another method that supplies that downstream result, not forced to materialize an unused array. Fairness does not mean every algorithm must use the same internal representation.
\head{How the original measurements are treated here}
The preserved manuscript contains its historical CPU/GPU corpus. Replaying its figure generator redraws those saved observations. It does not rerun the recorded device, establish present hardware performance, or transform logical avoided work into physically executed floating-point operations. No new GPU timing is added by this appendix.

Exact representation gains, device gains, kernel gains, and amortization gains can be isolated only through appropriately matched rows. A total ratio that mixes them is still an observation, but not a decomposition theorem.
\takeaway{A timing number without its case and scope is incomplete data. Preserve the mathematical output contract before claiming a speedup.}

\vpage{27}{Close the loop on the running example}{FACTORS, APPLICATION, AND AN INDEPENDENT REFERENCE}{Main text: Theorems 2.3, 3.4, and 5.1, pp. 4, 7, and 13.}
Return to the two operators from \pagerefV{9}. The exact product representation has diagonal part $D_AD_B$ and two-column correction factors $U_C,V_C$. Apply it to the original $x=(1,-1,2,0,-2,3)^T$.
\[
V_C^Tx=(2,13)^T,\qquad
(AB)x=D_AD_Bx+U_C(V_C^Tx).
\]
\vfig{end_to_end}
\head{A complete reader calculation}
The local contribution is $(1,-5,18,6,-12,44)^T$. The correction is $(25,21,25,21,25,21)^T$. Their sum is
\[
(26,16,43,27,13,65)^T.
\]
The checker forms the dense matrices $A$ and $B$ independently, multiplies them, applies the product to $x$, and compares all six components. The dense array is a validation reference, not a requirement of the factored application route.
\head{Choose what happens next}
For another input, reuse the factors. For a required explicit matrix, pay materialization. For a long product chain, track rank growth and decide whether recompression is useful. If approximation is introduced, transport its error through the later operators rather than reporting only a local singular tail.

This small example does not predict a speedup: at dimension six, implementation overhead can dominate. Its purpose is to make the semantics auditable. The same identity can then be used inside a larger, separately benchmarked workload.
\takeaway{The full route is: represent the exact operator, retain every interaction, return the requested object, and validate against an independent reference. Performance is a further experiment.}

\vpage{28}{What this edition certifies}{FORMAL, EXACT, NUMERICAL, AND EMPIRICAL EVIDENCE}{Reading source: sixth-block supplement, pp. 1 and 7--8; current release ledger.}
The preserved v0.6.0 release records 56 of 56 manuscript-proved finite items, 392 named theorem declarations in 40 modules, and 66 separately counted construction audits. Of the theorem declarations, 19 are reused foundation results. These are inherited records, not a new Lean run for this editorial edition.

\begin{tabularx}{\linewidth}{@{}p{.25\linewidth}Y@{}}\toprule
Evidence & Its precise use here \\\midrule
Saved Lean records & Preserve the exact hypotheses and declared finite formalization boundary.\\
New exact checks & Reproduce the appendix's matrices, identities, integer models, and finite optima. They are not extracted Lean proofs.\\
Vector drawings & Display explicit arrays, schedules, and sampled geometric formulas. Pixel accuracy is not algebraic certification.\\
Replayed Python tests & Exercise the preserved software on its stated finite test cases. Floating-point tests retain numerical tolerances.\\
Historical GPU corpus & Remains a record of the supplied original cases and environment. It is not a new measurement.\\\bottomrule
\end{tabularx}
\head{The recompression distinction}
The sixth block derives best-rank and singular-tail conclusions from ordered thin-SVD equations, rather than assuming an optimality certificate. It does not independently establish general SVD existence, certify a floating-point QR/SVD algorithm, or identify every result with a canonical singular-values API. Those are separate statements.
\head{The four classical entries remain visible}
The manuscript separately tracks Eckart--Young--Mirsky, the conjugate-gradient energy estimate, the floating-point product model, and simultaneous spectral decomposition. The strengthened ordered-SVD proof does not silently promote all associated existence and numerical-algorithm claims.
\head{Limits of exact compilation}
The optimizer needs a finite candidate tree with additive data. Nonlocal cache effects, overlapping decompositions, global rank coupling, and unrestricted discovery can destroy its optimal substructure. No new generic dense-matrix exponent or universal accelerator claim follows from this appendix.
\takeaway{This is a reader-facing reconstruction of the paper's finite concepts, with reproducible examples and preserved source evidence. It changes exposition, not the formal theorem contract.}
